\lecture{7}{28. Oktober 2025}{Fri svingning/egensvingning af MDOF-systemer (del II)}

\begin{problemwithimage}{p7_1}{7.1}
  Vi returnerer til systemet behandlet i opgave 6.1 og tilføjer en dæmpningsmekanisme som vist.
  Bevis, at dæmpningen er klassisk.
\end{problemwithimage}
  
\begin{figure}[ht]
    \centering
    \incfig[0.9]{p7_1fbd}
    \caption{Free body diagram of the system}
    \label{fig:p7_1fbd}
\end{figure}

\begin{derivation}
  Baseret på \Cref{fig:p7_1fbd} kan Newtons 2. lov anvendes som
  \begin{align*}
    m \cdot \ddot{d}_1 &= 2k d_2 - 3k d_1 + c (\dot{d}_2 - \dot{d}_1) \\
    m \cdot \ddot{d}_1 + 3k d_1 - 2k d_2 + c \dot{d}_1 - c \dot{d}_2 &= 0 \\
    m \cdot \ddot{d}_2 &= - 3k d_2 + 2kd_1 - c (\dot{d}_2 - \dot{d}_1) \\
    m \cdot \ddot{d}_2 - 2k d_1 + 3k d_2 + c \dot{d}_2 - c \dot{d}_1 &= 0
  .\end{align*}
  Dermed har vi
\[ 
  M = \begin{bmatrix}
  m & 0\\
  0 & m\\
  \end{bmatrix}, \quad C = \begin{bmatrix}
  c & -c\\
  -c & c\\
  \end{bmatrix}, \quad K = \begin{bmatrix}
  3k & -2k\\
  -2k & 3k\\
  \end{bmatrix}
.\]
Fra Opgave 6.1 har vi at $\Phi = \begin{bmatrix}
1 & 1\\
1 & -1\\
\end{bmatrix}$
  Desuden har vi $\tilde{C}$ givet som
  \[ 
    \tilde{C} \equiv \Phi^{\intercal} C \Phi 
  .\]
  Dermed er $\tilde{C}$:
  \[
    \tilde{C} = \begin{bmatrix}
    1 & 1\\
    1 & -1\\
    \end{bmatrix}^{\intercal} \begin{bmatrix}
    c & -c\\
    -c & c\\
    \end{bmatrix} \begin{bmatrix}
    1 & 1\\
    1 & -1\\
    \end{bmatrix} = \begin{bmatrix}
    0 & 0\\
    0 & 4c\\
    \end{bmatrix}
  .\]
  Eftersom $\tilde{C} = \begin{bmatrix}
  0 & 0\\
  0 & 4c\\
  \end{bmatrix}$ er diagonal må der derfor være tale om klassisk dæmpning.
\end{derivation}

\begin{problemwithimage}{p7_2}{7.2}
  Vi betragter det viste rammesystem.
  Heri gælder det, at de horisontale bjælker er uendeligt stive og har masserne $m_1$ og $m_2$, mens de vertikale bjælker (også kaldet søjler) er fleksible (alle med stivheden $k_1$) og masseløse.
  En fjeder med stivheden $k_2$ er påsat den øverste horisontale bjælke.
  Systemet modelleres som et 2DOF-system, hvori den første frihedsgrad, $d_1(t)$, er den horisontale flytning af den nederste horisontale bjælke, mens den anden frihedsgrad, $d_2(t)$, er den horisontale flytning af den øverste horisontale bjælke.
  Det oplyses, at de udæmpede vinkelegenfrekvenser og de tilhørende masse-normaliserede egenvektorer af 2DOF-systemet er
  \[ 
  \omega_1 = \sqrt{200} \, \unit{rad\per\s}, \qquad \omega_2 = \sqrt{600} \, \unit{rad\per\s}, \qquad \phi_1 = \begin{bmatrix}
  1 \\
  1 \\
  \end{bmatrix}, \qquad \phi_2 = \begin{pmatrix}
  1\\
  -1\\
  \end{pmatrix}
  .\]
  a) Bestem værdierne for $k_1$ og $k_2$.

  b) Bestem værdierne for $m_1$ og $m_2$.
  
  Hvad angår 2DOF-systemet, så oplyses begyndelsesbetingelserne $d_0 = \begin{pmatrix}
  \num{0,1}  & \num{0,2} \\
  \end{pmatrix}^{\intercal}$ og $\dot{d}_0 = \begin{pmatrix}
  0 & 0\\
  \end{pmatrix}^{\intercal}$.
  Endvidere oplyses det, at systemet er klassisk dæmpet med de modale dæmpningskonstanter $\tilde{c}_{11} = 8$ og $\tilde{c}_{22} = 10$.

  c) Bestem dæmpningsmatricen, $C \in \mathbb{R}^{2 \times 2}$.

  d) Bestem begyndelsesaccelerationsvektoren, $\ddot{d}_0$.
\end{problemwithimage}

\begin{derivation}
  Vi har fået givet to egenfrekvenser og de dertilhørende egenvektorer.
  Vi har altså:
  \[ 
  \Lambda = \begin{bmatrix}
  \qty{200}{\s^{-1}}  & 0\\
  0 & \qty{600}{\s^{-1}}\\
  \end{bmatrix}, \qquad \Phi = \begin{bmatrix}
  1 & 1\\
  1 & -1\\
  \end{bmatrix}, \qquad \tilde{M} = I
  .\]
  Fra undervisningen har vi at siden vi generelt har
  \[ 
  \Phi^{\intercal} K \Phi = \Phi^{\intercal} M \Phi \Lambda
  \]
  så medfører $\tilde{M}= I$ at $\tilde{K} = \Lambda$.
  $\tilde{K}$ er givet som:
  \[ 
  \tilde{K} \equiv \Phi^{\intercal} K \Phi
  .\]
  Altså:
  \begin{align*}
    \tilde{K} &= \Phi^{\intercal} K \Phi \\
    K &= \left( \Phi^{\intercal} \right)^{-1} \tilde{K} \Phi^{-1} \\
      &= \begin{bmatrix}
      200 & -100\\
      -100 & 200\\
      \end{bmatrix}
  .\end{align*}
  Vi har
  \[ 
  K = \begin{bmatrix}
  4k_1 & -2k_1\\
  -2k_1 & 2k_1 + k_2\\
  \end{bmatrix} \quad \implies \quad k_1 = 50, \: k_2 = 100
  .\]
  Samme procedure kan følges for masserne.
  Siden $\tilde{M} = I$ fås:
  \begin{align*}
    \tilde{M} &= \Phi^{\intercal} M \Phi \\
    M &= \left( \Phi^{\intercal} \right)^{-1} \Phi \\
      &= \begin{bmatrix}
      \frac{1}{2} & 0\\
      0 & \frac{1}{2}\\
      \end{bmatrix}
  .\end{align*}
  Dermed har vi $m_1 = m_2 = \frac{1}{2}$.

  Udnyttes samme teknik for dæmpningsmatricen får vi
  \begin{align*}
    \tilde{C} &= \Phi^{\intercal} C \Phi \\
    C &= \left( \Phi^{\intercal} \right)^{-1} \tilde{C} \Phi^{-1} \\
      &= \begin{bmatrix}
      \frac{1}{2} & \frac{1}{2}\\
      \frac{1}{2} & - \frac{1}{2}\\
      \end{bmatrix} \begin{bmatrix}
      8 & 0\\
      0 & 10\\
      \end{bmatrix} \begin{bmatrix}
      \frac{1}{2} & \frac{1}{2}\\
      \frac{1}{2} & -\frac{1}{2}\\
      \end{bmatrix} \\
      &= \frac{1}{2} \begin{bmatrix}
      9 & -1\\
      -1 & 9\\
      \end{bmatrix}
  .\end{align*}
  
  For at bestemme $\ddot{d}_0$ benytter vi den generelle formel
  \begin{align*}
    M \ddot{d}_0 + C \dot{d}_0 + K d_0 &= 0 \\
    \ddot{d}_0 &= - M^{-1} \left( C \dot{d}_0 + K d_0 \right) \\
     &= - \begin{bmatrix}
     \frac{1}{2} & 0\\
     0 & \frac{1}{2}\\
     \end{bmatrix}^{-1} \left( \frac{1}{2} \begin{bmatrix}
     9 & -1\\
     -1 & 9\\
     \end{bmatrix} \cdot \begin{bmatrix}
     0 \\
     0 \\
     \end{bmatrix} + \begin{bmatrix}
     200 & -100\\
     -100 & 200\\
     \end{bmatrix} \cdot \begin{bmatrix}
     \num{0,1} \\
     \num{0,2} \\
     \end{bmatrix}\right) \\
     &= \begin{bmatrix}
     0\\
     -60\\
     \end{bmatrix}
  .\end{align*}
  
\end{derivation}

\finalresult{
\begin{aligned}
    \mathrm{a}) &\quad k_1 = \num{50}, \, k_2 = 100 \\ 
    \mathrm{b}) &\quad m_1 = m_2 = \frac{1}{2} \\
    \mathrm{c}) &\quad C = \frac{1}{2} \begin{bmatrix}
    9 & -1 \\
    -1 & 9\\
    \end{bmatrix} \\
    \mathrm{d}) &\quad \ddot{d}_0 = \begin{bmatrix}
    0\\
    -60\\
    \end{bmatrix}
.\end{aligned}
}
